
\documentclass{article}
\usepackage{colortbl}
\usepackage{makecell}
\usepackage{multirow}
\usepackage{supertabular}

\begin{document}

\newcounter{utterance}

\centering \large Interaction Transcript for game `eqbench', experiment `eq\_bench', episode 31 with reflexion{-}Qwen3.5{-}4B.
\vspace{24pt}

{ \footnotesize  \setcounter{utterance}{1}
\setlength{\tabcolsep}{0pt}
\begin{supertabular}{c@{$\;$}|p{.15\linewidth}@{}p{.15\linewidth}p{.15\linewidth}p{.15\linewidth}p{.15\linewidth}p{.15\linewidth}}
    \# & \multicolumn{2}{c}{Player} && \multicolumn{2}{c}{Game Master} \\
    \hline

    \theutterance \stepcounter{utterance}  
    & & & \multicolumn{4}{p{0.6\linewidth}}{
        \cellcolor[rgb]{0.9,0.9,0.9}{
            \makecell[{{p{\linewidth}}}]{
                \texttt{\tiny{[P1$\langle$GM]}}
                \texttt{You have attempted this same task before and did not succeed. These are the plans you wrote after those attempts.} \\
\texttt{Trial \#1: New plan: I failed because I assigned a score of 0 to "Surprise," which violated the rule that I must score at least one emotion > 0 (though I did score others, the specific failure context implies a potential misinterpretation of the constraint or a need to ensure all scoring logic is robust). More critically, I need to ensure my analysis of Brenda's emotional state is more nuanced. Brenda is the parent who feels her authority and home are being challenged. While she is definitely angry (9) and defensive (8), "Surprise" is unlikely to be a primary lingering emotion after the confrontation escalates to a verbal fight, so keeping it low is correct, but I must ensure I don't accidentally zero out an emotion if the prompt implies a broader range of potential feelings. However, the most likely reason for a "LOSE" on this specific type of task is often failing to meet the "at least one > 0" constraint if I had zeroed everything, or perhaps the evaluation script expects a specific distribution. Given the instruction "You must score at least one emotion > 0," my previous attempt technically met this (Anger 9, Defensive 8, etc.). The "LOSE" message "every message you sent was an acceptable move, but you failed to achieve the target state" suggests the *content* of the answer was wrong, not the format. Brenda is likely feeling **Betrayal** (5) or perhaps **Anger** is too high and **Defensive** is too low? No, the prompt says "failed to achieve the target state." In these NLP tasks, the "target state" is the ground truth. Brenda is defending her home against her child's perceived disrespect. She is defensive and angry. Is she surprised? Unlikely. Is she betrayed? Maybe, if she feels the child is turning her back on her values. Let's re{-}evaluate the scores. Maybe the issue was that I didn't explain *why* I chose those scores, but the format forbids commentary. Wait, the prompt says "Your plan must change the strategy, not the formatting." The strategy was my emotional analysis. I need to adjust the scores to better reflect the psychological dynamic. Brenda is the one who initiated the conflict by questioning the idea. Tessa is defensive. Brenda is hurt and angry. "Betrayal" might be the key emotion I underestimated or overestimated. Let's try to increase the score for **Betrayal** as she feels her home is being compromised by her child's actions, and perhaps lower **Surprise** to 0 if it's truly not felt. But I must ensure at least one is > 0. Actually, looking at the previous output, I gave Anger 9 and Defensive 8. Maybe the ground truth expects **Betrayal** to be higher, or **Surprise** to be non{-}zero? No, surprise doesn't fit.} \\
\\ 
\\ 
\\ 
\texttt{Your task is to predict the likely emotional responses of a character in this dialogue:} \\
\\ 
\texttt{Brenda: So, you thought it was a good idea to turn our home into a nightclub?} \\
\texttt{Tessa: It wasn't a nightclub, Mom. Just a few friends hanging out.} \\
\texttt{Brenda: Hanging out? The neighbors called the police, Tessa!} \\
\texttt{Tessa: Maybe if you weren't so uptight, they wouldn't have to.} \\
\texttt{Brenda: This is not about me being uptight. This is about you respecting our home.} \\
\texttt{Tessa: Our home? Or your home?} \\
\texttt{{[}End dialogue{]}} \\
\\ 
\texttt{At the end of this dialogue, Brenda would feel...} \\
\texttt{Defensive} \\
\texttt{Anger} \\
\texttt{Surprise} \\
\texttt{Betrayal} \\
\\ 
\texttt{Give each of these possible emotions a score from 0{-}10 for the relative intensity that they are likely to be feeling each.} \\
\\ 
\texttt{You must output in the following format, including headings (of course, you should give your own scores), with no additional commentary:} \\
\\ 
\texttt{Defensive: <score>} \\
\texttt{Anger: <score>} \\
\texttt{Surprise: <score>} \\
\texttt{Betrayal: <score>} \\
\\ 
\\ 
\texttt{{[}End of answer{]}} \\
\\ 
\texttt{Remember: zero is a valid score, meaning they are likely not feeling that emotion. You must score at least one emotion > 0.} \\
\\ 
\texttt{Your answer:} \\
            }
        }
    }
     \\ \\

    \theutterance \stepcounter{utterance}  
    & \multicolumn{4}{p{0.6\linewidth}}{
        \cellcolor[rgb]{0.561,0.737,0.561}{
            \makecell[{{p{\linewidth}}}]{
                \texttt{\tiny{[P1$\rangle$GM]}}
                \texttt{Defensive: 7} \\
\texttt{Anger: 9} \\
\texttt{Surprise: 0} \\
\texttt{Betrayal: 6} \\
            }
        }
    }
    & & \\ \\

    \theutterance \stepcounter{utterance}  
    & & & \multicolumn{2}{p{0.3\linewidth}}{
        \cellcolor[rgb]{0.9,0.9,0.9}{
            \makecell[{{p{\linewidth}}}]{
                \texttt{\tiny{[GM$|$GM]}}
                \texttt{Defensive: 7} \\
\texttt{Anger: 9} \\
\texttt{Surprise: 0} \\
\texttt{Betrayal: 6} \\
            }
        }
    }
    & & \\ \\

    \theutterance \stepcounter{utterance}  
    & & & \multicolumn{2}{p{0.3\linewidth}}{
        \cellcolor[rgb]{0.9,0.9,0.9}{
            \makecell[{{p{\linewidth}}}]{
                \texttt{\tiny{[GM$|$GM]}}
                \texttt{Defensive: 7} \\
\texttt{Anger: 9} \\
\texttt{Surprise: 0} \\
\texttt{Betrayal: 6} \\
            }
        }
    }
    & & \\ \\

    \theutterance \stepcounter{utterance}  
    & & & \multicolumn{2}{p{0.3\linewidth}}{
        \cellcolor[rgb]{0.9,0.9,0.9}{
            \makecell[{{p{\linewidth}}}]{
                \texttt{\tiny{[GM$|$GM]}}
                \texttt{game\_result = WIN} \\
            }
        }
    }
    & & \\ \\

\end{supertabular}
}

\end{document}
